\documentclass{article}
\usepackage{graphicx}
\usepackage{amsmath}
\usepackage{booktabs}
\usepackage{array}
\usepackage{enumitem}
\usepackage{xcolor}
\usepackage{titlesec}
\usepackage{fancyhdr}
\usepackage{listings}

\title{Architectural Styles}
\author{0xSaad}
\date{\today}

\begin{document}

\maketitle

There are different ways to view the organization of distributed systems:

\begin{itemize}
    \item Software Architecture: The logical organization.
    \item System Architecture: The actual physical realization.
\end{itemize}

\subsection{Component}

A modular unit with well-defined interfaces that is replaceable within its environment.

\subsection{Connector}

A mechanism that mediates communication, coordination, or cooperation among components.

Examples:
\begin{itemize}
    \item RPC (Remote Procedure Call)
    \item Messaging
    \item Streaming
\end{itemize}

\subsubsection{Event-Based Architectures}

Components communicate via events using the Publish/Subscribe paradigm.

\begin{itemize}
    \item Components can publish and subscribe to events.
    \item Events are delivered to subscribers.
    \item Processes are decoupled in space.
    \item Processes are coupled in time.
\end{itemize}

Types:

\begin{itemize}
    \item Subject-Based: A message belongs to groups, channels, or topics. Users subscribe to these groups and receive associated messages.
    \item Content-Based: Message delivery is based on a query or predicate issued by the subscriber.
\end{itemize}

\subsubsection{Object-Based Architectures}

\begin{itemize}
    \item Components interact via RPC.
    \item Each object corresponds to a software component.
\end{itemize}

A remote object provides a service comprising a set of functions.

The service description includes function names, argument numbers and types, and return types.

\paragraph{CORBA}

Common Object Request Broker Architecture.

An industry standard for distributed objects programming.

\paragraph{Java RMI}

Remote Method Invocation allows an object running in one Java Virtual Machine to invoke methods on an object running in another Java VM.

\paragraph{Amazon S3}

Objects (files) are placed into buckets (directories).

Buckets cannot be placed inside other buckets.

\begin{itemize}
    \item PUT: Create a bucket or object.
    \item GET: List objects in a bucket.
    \item GET: Read an object using its URI.
\end{itemize}

\section{What Are the Different Software Architectures?}

\begin{itemize}
    \item Layered Architectures
    \item Event-Based Architectures
    \item Data-Centered Architectures
    \item Object-Based Architectures
    \item Shared Data-Space Architectures
\end{itemize}

\section{System Architectures}

\begin{itemize}
    \item Centralized: Client-Server systems and Multitiered systems (Servers as clients).
    \item Decentralized: Peer-to-Peer systems.
    \item Hybrid: Client-Server + Decentralized.
\end{itemize}

\subsection{Layered Architecture}

Components are organized in a layered fashion.

A component at layer $L_i$ is allowed to call components at the lower layer $L_{i-1}$, not the other way around.

Control flows from layer to layer (down), and responses flow upwards (up).

\subsubsection{Example: OSI Model}

The OSI model is a conceptual model that characterizes and standardizes communication functions.

It consists of seven layers:

\begin{enumerate}
    \item Application
    \item Presentation
    \item Session
    \item Transport
    \item Network
    \item Data Link
    \item Physical
\end{enumerate}

\IfFileExists{images/osi.png}{
    \includegraphics[width=0.5\textwidth]{images/osi.png}
}{
    \textit{OSI diagram: images/osi.png}
}

\subsubsection{Levels}

\begin{itemize}
    \item User-Interface Level: Provides access to the user with the system.
    \item Processing Level: Contains the applications in the system.
    \item Data Level: Manages the data in the system.
\end{itemize}

\subsection{Data-Centered Architecture}

Processes communicate through a common repository.

\begin{itemize}
    \item Evolves around the data.
    \item Communication takes place through files.
    \item Examples: Web-based distributed systems and Distributed Database Systems.
\end{itemize}

\subsection{Shared Data-Space Architecture}

Combines Event-Based and Data-Centric architectural styles.

Processes communicate through a common repository.

\begin{itemize}
    \item Publishers post items to a shared data space.
    \item Subscribers access items at a later time.
    \item Processes are decoupled in both space and time.
\end{itemize}

\subsubsection{JavaSpaces}

A JavaSpace is a Jini service that stores Java objects in memory.

It supports collaboration among services.

Operations:

\begin{itemize}
    \item Write: Write an object to JavaSpace.
    \item Read: Return a copy and leave the object in JavaSpace.
    \item Take: Return a copy and remove the object from JavaSpace.
    \item Notify: Register a callback to be notified when changes occur.
\end{itemize}

\subsection{Client-Server Architecture}

\begin{itemize}
    \item Client: A process that requests a service from a server.
    \item Server: A process that offers a service, such as a file or database service.
\end{itemize}

\subsubsection{Steps at Client}

\begin{enumerate}
    \item Make a message for the server.
    \item Identify the service the client wants.
    \item Send the message to the server.
\end{enumerate}

\subsubsection{Steps at Server}

\begin{enumerate}
    \item Wait for an incoming request.
    \item Process the request.
    \item Package the result in a reply message.
    \item Send the reply to the client.
\end{enumerate}

Communication between client and server can use:

\begin{itemize}
    \item Connectionless Protocol (UDP): Used when the underlying network is reliable, according to the slides.
    \item Connection-Oriented Protocol (TCP): Used when the underlying network is unreliable, according to the slides.
\end{itemize}

\subsubsection{UDP Problems}

\begin{itemize}
    \item Occasional transmission failures may occur.
    \item The client resends the request when no reply arrives.
    \item The client cannot determine whether the request or reply was lost.
    \item Resending a request may execute the operation twice.
\end{itemize}

\subsubsection{TCP Steps}

\begin{enumerate}
    \item The client establishes a connection with the server.
    \item The client sends the request.
    \item The server sends the result using the same connection.
    \item The connection is closed after receiving the result.
\end{enumerate}

Disadvantages mentioned in the slides:

\begin{itemize}
    \item Low performance.
    \item Connection setup is costly for small messages.
\end{itemize}

\subsection{Multitiered Architecture}

Application programs are distributed across different machines.

\subsubsection{Two-Tier Organization}

\begin{itemize}
    \item Client Machine: Contains programs implementing the user-interface level.
    \item Server Machine: Contains the processing and data levels.
\end{itemize}

\subsubsection{Three-Tier Organization}

\begin{itemize}
    \item User Interface: Presentation level.
    \item Application Server: Processing level.
    \item Database Server: Data level.
\end{itemize}

Request flow:

\begin{enumerate}
    \item User interface requests an operation.
    \item Application server requests data.
    \item Database server returns data.
    \item Application server returns the result.
\end{enumerate}

\subsection{Layers vs. Tiers}

\begin{itemize}
    \item Layer: Logical layer, a term in Software Architecture.
    \item Tier: Physical layer, a term in System Architecture.
\end{itemize}

\subsection{P2P Systems}

Peer-to-Peer networks use distribution between equivalent peers.

This is called horizontal distribution, as opposed to layered vertical distribution.

Problems:

\begin{itemize}
    \item Who should store a particular data item?
    \item How is a message routed through the peer group?
    \item Sending messages to all peers may flood the network.
\end{itemize}

\subsection{Overlay Network}

An overlay network is a logical network built on top of a physical network.

Focus at the application level.

\subsection{Hybrid Architecture}

Hybrid Architecture combines Decentralized and Client-Server systems.

\subsection{Edge-Server Systems}

Servers are placed at the edge of the network.

They are located at the boundary between enterprise networks and the Internet, such as an ISP.

\subsection{Collaborative Distributed Systems}

Components of a collaborative content distribution network:

\begin{itemize}
    \item A component that redirects client requests to other servers.
    \item A component for analyzing access patterns.
    \item A component for managing the replication of Web pages.
\end{itemize}

\end{document}